\project{06}{Sampling on a timer at exactly 1 kHz}{NUCLEO-H7A3ZI-Q with an MCC 118 on a Raspberry Pi}{Timer-triggered acquisition, proving the rate externally}

\begin{keyfacts}
\item[Board] NUCLEO-H7A3ZI-Q, plus a Raspberry Pi 4 carrying an MCC 118 DAQ HAT as the external witness
\item[Peripherals] TIM6 as the trigger, ADC1 with hardware trigger, DMA in double-buffer mode, TIM2 and TIM5 as instruments, one free GPIO as the marker
\item[Toolchain] arm-none-eabi-gcc with CMake as in chapter 1; Python 3 with the HAT library on the Pi
\item[Operating system] Bare metal on the board, Raspberry Pi OS on the witness
\item[Difficulty] 4 of 5
\item[Effort] 4 evenings of about four hours, one of them entirely on the Pi
\item[Deliverable] Three sampling builds that share one application, an external measurement that states the sample rate with an uncertainty, and the two instruments the rest of this book measures with
\end{keyfacts}

\subsection*{Why this project}

Every chapter after this one makes a claim about time. Chapter 7 claims a wake
takes so many milliseconds, chapter 10 splits a duty cycle into phases, chapter
18 counts cycles per transform and chapter 20 reports a latency distribution.
None of those claims mean anything until there is an instrument, and this bench
has no oscilloscope and no logic analyser. So this chapter does two jobs at
once. It builds a sampling front end that acquires at a fixed rate while the
processor is busy elsewhere, and it builds the instruments: a cycle counter
wrapper with a fallback, and a gated frequency counter that can be pointed at
any periodic signal on the board.

The second job is the reason the chapter is worth four evenings rather than
one. A sample rate that is configured is not a sample rate that is proven. The
usual practice is to write the prescaler and the reload value, do the
arithmetic, and call the result 1 kHz. That arithmetic is worth doing and this
chapter does it, but it establishes only that the divider chain is exact. It
says nothing about the oscillator the chain divides, nothing about whether the
interrupt actually fires when the timer says it does, and nothing about jitter.
The only honest way to settle those is to put a second, independent clock next
to the first and compare them. The MCC 118 on a Raspberry Pi is that second
clock: it samples a pin the board toggles from its own sampling interrupt, at a
rate high enough that each edge lands between two known instants.

The third reason is structural. A sampling loop that does its work inside the
interrupt looks correct at 1 kHz and fails silently at 10 kHz, because the
overrun has nowhere to be reported. Decoupling acquisition from processing is
the pattern the rest of Part I and all of Part II depend on, and this chapter
builds it three ways and measures the difference.

\begin{note}{This part is not the STM32H7 the internet is about}
The timer inventory, the transfer request numbers that connect the converter to
a transfer stream, and the rule about which transfer engine reaches which memory
all differ between RM0455, which documents this part, and RM0433, which
documents its better-known sibling. The prior art named below was written on a
Cortex-M7 Nucleo of the other kind. Its structure transfers; its request
numbers do not. Read them out of RM0455's request table for this part, not out
of a working project for another one.
\end{note}

\subsection*{Prior art and what to reuse}

\begin{table}[H]\centering\small
\begin{tabular}{p{32mm}p{46mm}p{40mm}p{20mm}}
\toprule
Source & What it gives & What it does not & Licence \\
\midrule
Ian Ross's five-step progression on a Cortex-M7 Nucleo & Exactly this chapter's shape: the same acquisition solved in rising order of sophistication, with a written series explaining each step's reason & It is small, unmaintained, and written for a different member of this family. No double buffering, no external verification, no instrument & MIT \\
The vendor's basic examples for this board & Working register settings for a timer trigger into the converter, in a tree that also solves the same job polled and on interrupt & Its example projects sit under a licence tied to this vendor's silicon; only the basic examples are permissive. No verification of the resulting rate & BSD-3-Clause for the basic examples \\
Majerle's transfer-driven serial work & The half-transfer and transfer-complete pattern this chapter reuses for the converter, proven across eight families & It is a serial receive problem, not acquisition, and it does no cache maintenance & MIT \\
GorgonMeducer's performance counter & A cycle-accounting library that deliberately runs on the system tick rather than the core's trace counter, because the tick is universally present & It is a whole library where this chapter needs a twenty-line wrapper with a self-test & Apache-2.0 \\
The HAT library for the acquisition board & A Python interface to the external witness with a scan API and published examples & Its licence is not recorded in this volume's prior-art pool. Read the repository's licence file before vendoring any of it. This chapter only calls it & Read before use \\
\bottomrule
\end{tabular}
\caption{Prior art for chapter 6. The five-step progression is the closest thing to this chapter that exists, and its written series is worth reading before you start; the part it is written for is not this one.}
\end{table}

What is left to write is the verification and the instruments. No source in the
pool measures the rate it configures. The external witness method below, the
edge-interpolation arithmetic that turns a sampled square wave into a rate with
an uncertainty, the cycle counter wrapper with a runtime self-test and a timer
fallback, and the gated counter and its reciprocal upgrade are all the author's
work, and they are the reason this chapter comes before the seven that cite it.

\subsection*{Parts from the inventory}

\begin{table}[H]\centering\small
\begin{tabular}{p{40mm}p{66mm}p{40mm}}
\toprule
Part & Role & Interface \\
\midrule
NUCLEO-H7A3ZI-Q & The subject. Samples on a timer and toggles a marker pin from the sampling interrupt & Micro USB to the host \\
Raspberry Pi 4 & The witness. Runs the scan and the analysis & Ethernet or wireless to the host \\
MCC 118 DAQ HAT & The calibrated reference, 8 channels, 12-bit, 100 kS/s aggregate & 40-pin header on the Pi, screw terminals to the board \\
Two jumper wires & Marker pin and a common ground reference & Screw terminal at one end, Zio header at the other \\
A 220 \textohm{} resistor & In series with the marker pin, so a wiring error costs nothing & Breadboard \\
\bottomrule
\end{tabular}
\caption{Inventory items used in chapter 6. Nothing is bought and nothing is soldered. Whether male-to-male jumper wires exist in the drawer is item 36 on the confirm list; the screw terminals on the HAT accept a bare wire end either way.}
\end{table}

\subsection*{System architecture}

\diagram{c06_arch}{Two independent clocks, one signal between them. The board's divider chain produces the sample instants; the witness has its own oscillator and never sees the board's. What the comparison proves is agreement, which is a stronger and more modest claim than accuracy.}

The board's side is a chain with no software in it: a timer reaches its reload
value, emits a trigger event, the converter starts on that event, and the
transfer engine moves the result to memory. The processor is told only when a
half of the buffer is complete, and between those notifications it is free. The
witness is deliberately dumb: it samples one channel continuously at a high
rate, writes the samples to a file with a start timestamp, and the analysis
happens afterwards on the same Pi. Nothing on it is synchronised to the board,
which is the point.

\subsection*{Peripheral configuration}

\begin{table}[H]\centering\small
\begin{tabular}{p{20mm}p{28mm}p{24mm}p{30mm}p{30mm}}
\toprule
Peripheral & Mode & Clock source & Pins and function & Interrupt and transfers \\
\midrule
TIM6 & Up counter, update event on TRGO & APB1 timer clock & None & No interrupt in the timer build \\
ADC1 & Single channel, hardware trigger on the timer's TRGO & ADC kernel clock, confirm the source in RM0455 & One analogue input, channel confirmed in the datasheet & Transfer request to a DMA stream \\
DMA1 stream & Double buffer, two memory pointers & Bus clock & None & Half and full transfer interrupts \\
GPIO marker & Push-pull output, toggled in the sampling interrupt & Peripheral bus & A free Zio pin, item 13 on the confirm list & None \\
TIM2 & Free running 32-bit, or input capture for the reciprocal counter & APB1 timer clock & Capture pin from the alternate-function table & Capture interrupt, reciprocal build only \\
TIM5 and TIM7 & External clock mode 2 counting edges, and an up counter as the gate & The signal under test, and the APB1 timer clock & External trigger pin from the alternate-function table & Gate interrupt only \\
SysTick & 1 kHz reload, the simplest build & Core clock & None & SysTick exception \\
\bottomrule
\end{tabular}
\caption{Peripheral configuration. Two 32-bit timers exist on this part and both are spent on instruments rather than on the application, which is a deliberate allocation and is worth stating in the repository README.}
\end{table}

\subsection*{Wiring}

\diagram{c06_wiring}{One signal wire and one ground wire. The series resistor is not protection for the board, which drives the pin, but insurance against the wire finding the wrong terminal.}

Two safety points before the wires go in. The board's logic is 3.3 V and
nothing from the Pi or the HAT may be driven back into a board pin; the marker
is an output on the board and an input on the HAT, and it stays that way. The
grounds must be common or the witness measures an undefined potential, so the
ground wire is not optional and goes in first. The HAT's input range and
whether it offers an external trigger are item 35 on the confirm list: read the
HAT's user guide before assuming a 3.3 V swing is comfortably inside the range,
even though every range this HAT is likely to offer contains it.

\subsection*{Memory and timing budget}

\diagram{c06_mem}{Where 1 kHz actually comes from. The divider chain is exact, and that is a different claim from accurate. Everything to the left of the timer is inherited from chapter 1's clock configuration and its accuracy is the probe's.}

The arithmetic is worth writing out because it is the part that is exact. With
the peripheral bus prescaler set so the timer kernel clock is 280 MHz, a
prescaler value of 279 divides by 280 and gives a 1 MHz counter. A reload value
of 999 counts 1000 of those and gives an update event every millisecond. Both
divisions are integer, so there is no residue and no slow drift from rounding.
The same is true of the tick: a reload of 279999 on a 280 MHz core clock is
exactly one millisecond and fits in the tick's 24 bits with room to spare.

What is not exact is the 8 MHz the whole chain starts from, which comes from the
probe and carries the probe's tolerance. That number is not in this volume and
is not guessed at here. It is what the external measurement below reports.

\begin{table}[H]\centering\small
\begin{tabular}{p{44mm}p{28mm}p{28mm}p{30mm}}
\toprule
Quantity & Budget & Measured & Margin \\
\midrule
Sample period & 1.000000 ms by construction & not measured & not measured \\
Interrupt to marker edge & under 2 \textmu{}s & not measured & not measured \\
Edge to edge jitter & under 1 \textmu{}s peak & not measured & not measured \\
Cycles in the sampling interrupt & under 2000 of 280000 & not measured & not measured \\
Buffer, 128 samples per half & 512 B in a region the stream can reach & 512 B by construction & none needed \\
Processing budget per half & 128 ms & not measured & not measured \\
\bottomrule
\end{tabular}
\caption{The budget for chapter 6. Every measured cell is filled by the end of the chapter, by the witness for the first three rows and by this chapter's own cycle counter for the fourth. Nothing is filled by arithmetic.}
\end{table}

\subsection*{Firmware design (UML)}

\diagram{c06_uml}{One half of the buffer filling, with nothing to do in the interrupt but set a flag. The processing runs on the half that is no longer being written, which is what makes the block size a budget rather than a deadline.}

The design has one invariant and it is worth stating before the code. At any
instant the transfer engine is writing one half of the buffer and the
application may read the other. The interrupt's only job is to record which
half just became readable, and the only failure mode is that the application
has not finished with a half by the time the engine comes back to it. That
failure is detectable: the interrupt can see that the previous half was never
consumed and increment an overrun count. A design where the overrun is not
countable is the design this chapter replaces.

\subsection*{Data flow (ASCII)}

\begin{asciiart}
  board                                                      witness
  +-----------------------------------------------+          +----------------------------+
  | TIM6 update -> TRGO                            |          | Raspberry Pi 4             |
  |   |                                            |          |   MCC 118 DAQ HAT          |
  |   v                                            |          |   1 channel, 100 kS/s      |
  | ADC1 conversion (no software in the path)      |          |                            |
  |   |                                            |          |   scan -> samples.npy      |
  |   v                                            |          |        |                   |
  | DMA1 stream, double buffer                     |          |        v                   |
  |   +--> half A [128] --+                        |          |   threshold crossings      |
  |   +--> half B [128] --+                        |          |   linear interpolation     |
  |            |                                   |          |        |                   |
  |            v  half/full transfer interrupt     |          |        v                   |
  |   flag = which half is readable                |          |   straight-line fit        |
  |            |                                   |          |   slope = sample period    |
  |            v                                   |          |   residual = jitter        |
  |   main loop: feature over the readable half    |          +----------------------------+
  |            |                                   |                       ^
  |            v                                   |   3.3 V marker        |
  |   GPIO marker toggled in the interrupt --------+-----------------------+
  |                                               (220 ohm in series, common ground)
  +-----------------------------------------------+
\end{asciiart}

\subsection*{Repository layout}

\begin{plaincode}
nucleo-h7a3-sampling/
  CMakeLists.txt                  # three build targets from one application
  cmake/arm-none-eabi.cmake       # unchanged from chapter 1
  src/main.c                      # the application, identical in all three
  src/acq_systick.c               # build 1: sample in the tick exception
  src/acq_timer.c                 # build 2: timer TRGO starts the converter
  src/acq_dma_double.c            # build 3: double-buffered transfer
  src/acq.h                       # the one interface all three satisfy
  src/marker.c                    # the GPIO the witness watches
  instr/cyccnt.c                  # cycle counter wrapper, DWT or timer
  instr/cyccnt.h
  instr/freqcount.c               # gated counter, and the reciprocal upgrade
  instr/freqcount.h
  witness/scan.py                 # runs on the Pi, writes samples and metadata
  witness/rate.py                 # edges, fit, rate with an uncertainty
  witness/test_rate.py            # the analysis tested on synthetic input
  docs/measurement.md             # the method, written before the first run
  README.md
\end{plaincode}

\subsection*{Steps}

\begin{steps}
\item \textbf{Write the interface all three builds satisfy.} Decide this first,
because it is what makes the comparison honest. Three implementations behind one
header, selected by the build, so the application source is byte identical
across them.
\begin{ccode}
/* acq.h: one interface, three implementations. */
#include <stdint.h>
#include <stddef.h>

typedef struct {
    const uint16_t *samples;   /* the half that is safe to read */
    size_t          count;     /* samples in that half */
    uint32_t        seq;       /* increments once per half, never wraps in a run */
} acq_block_t;

void     acq_start(void);
int      acq_take(acq_block_t *out);   /* 1 if a block is ready, else 0 */
uint32_t acq_overruns(void);           /* halves the application never read */
\end{ccode}

\item \textbf{Build the cycle counter wrapper before anything else.} It is the
cheapest instrument on the board and the rest of the book uses it. Whether the
core's trace counter runs with no debugger attached is item 25 on the confirm
list, so the wrapper does not assume: it enables the counter, checks that it
moved, and reports failure rather than returning zeros.
\begin{ccode}
/* instr/cyccnt.c */
#include <stdint.h>
#include "stm32h7xx.h"          /* CMSIS device header for this part */

/* The debug block's lock access register. Some Cortex-M7 implementations
   require this key before the trace registers accept writes; on others the
   register reads as zero and the write is harmless. Do not decide which
   applies by reading a forum post: the self-test below decides it. */
#define DWT_LAR (*(volatile uint32_t *) 0xE0001FB0u)
#define DWT_LAR_KEY 0xC5ACCE55u

static int cyccnt_ok;

int cyccnt_init(void)
{
    CoreDebug->DEMCR |= CoreDebug_DEMCR_TRCENA_Msk;
    DWT_LAR = DWT_LAR_KEY;
    DWT->CYCCNT = 0u;
    DWT->CTRL  |= DWT_CTRL_CYCCNTENA_Msk;

    uint32_t a = DWT->CYCCNT;
    for (volatile int i = 0; i < 64; i++) { }
    cyccnt_ok = (DWT->CYCCNT != a);
    return cyccnt_ok;                 /* 0 means: use the timer backend */
}

uint32_t cyccnt_now(void) { return DWT->CYCCNT; }
\end{ccode}
If the self-test returns zero, the fallback is a free-running 32-bit timer with
its prescaler set to one and its reload value left at the maximum, read through
the same header. It is one of only two 32-bit timers on this part, which is why
chapter 20 is careful about who owns it. The two backends differ in one way
that matters and the header says so: the trace counter counts core clocks and
the timer counts timer clocks, and on this part those are the same number only
because chapter 1 configured the buses that way. Store the tick rate next to
the reading rather than assuming it.

\item \textbf{Build the first acquisition: the system tick.} This is the
simplest thing that can work and it is a real variant, not a strawman. The tick
exception fires at 1 kHz, starts a conversion, waits for it, stores it, and
toggles the marker.
\begin{ccode}
/* acq_systick.c */
void acq_start(void)
{
    SysTick->LOAD = 280000u - 1u;     /* 280 MHz core clock, exactly 1 ms */
    SysTick->VAL  = 0u;
    SysTick->CTRL = SysTick_CTRL_CLKSOURCE_Msk | SysTick_CTRL_TICKINT_Msk
                  | SysTick_CTRL_ENABLE_Msk;
}

void SysTick_Handler(void)
{
    marker_toggle();                  /* first, so the witness sees the edge */
    ADC1->CR |= ADC_CR_ADSTART;
    while ((ADC1->ISR & ADC_ISR_EOC) == 0u) { }   /* the cost of this build */
    ring_put((uint16_t) ADC1->DR);
}

/* marker.c: one store, no read-modify-write, no branch. The witness can only
   see what this function shows it, so its own cost is part of the measurement.
   Measure it with the wrapper from step 2 and record the number: every later
   chapter that uses a marker pin inherits it. */
void marker_toggle(void)
{
    static uint32_t level;
    level ^= 1u;
    MARKER_PORT->BSRR = level ? (1u << MARKER_PIN)
                              : (1u << (MARKER_PIN + 16));
}
\end{ccode}
The waiting loop is the whole lesson. Measure it with the wrapper from step 2
and write the number in the budget table. At 1 kHz it is invisible; the point of
measuring it now is that chapter 13 samples faster.

\item \textbf{Build the second acquisition: a timer starts the converter.} Now
no software sits between the timer and the conversion. The timer's update event
is routed to its trigger output and the converter is configured to start on it.
\begin{ccode}
/* acq_timer.c */
void acq_start(void)
{
    RCC->APB1LENR |= RCC_APB1LENR_TIM6EN;
    TIM6->PSC = 280u - 1u;            /* timer clock 280 MHz -> 1 MHz */
    TIM6->ARR = 1000u - 1u;           /* 1 MHz -> 1000 Hz, exactly */
    TIM6->CR2 = (2u << TIM_CR2_MMS_Pos);          /* TRGO on update */
    TIM6->EGR = TIM_EGR_UG;

    /* ADC1: external trigger, rising edge, source number from RM0455's
       trigger table for THIS part. Do not copy it from an H743 project. */
    ADC1->CFGR = (ADC1->CFGR & ~ADC_CFGR_EXTSEL_Msk)
               | (ADC_TRIGGER_TIM6_TRGO << ADC_CFGR_EXTSEL_Pos)
               | (1u << ADC_CFGR_EXTEN_Pos);      /* rising edge */
    ADC1->CR  |= ADC_CR_ADSTART;
    TIM6->CR1  = TIM_CR1_CEN;
}
\end{ccode}
The marker now moves into the end-of-conversion interrupt, which is one step
further from the timer and so one step worse as a witness of the timer. Toggle
it in both places under a build switch and let the witness tell you how much
that step costs. That comparison is a measurement this book has not seen
published for this family.

\item \textbf{Build the third acquisition: double-buffered transfer.} The
stream controller on this family carries two memory pointers and a bit saying
which one it is currently filling. Set both pointers, enable double-buffer
mode, and take the half and full transfer interrupts.
\begin{ccode}
/* acq_dma_double.c */
#define HALF 128u
static uint16_t buf_a[HALF], buf_b[HALF];
static volatile uint32_t ready_seq, taken_seq;
static const uint16_t *ready_ptr;

void acq_start(void)
{
    DMA1_Stream0->PAR  = (uint32_t) &ADC1->DR;
    DMA1_Stream0->M0AR = (uint32_t) buf_a;
    DMA1_Stream0->M1AR = (uint32_t) buf_b;
    DMA1_Stream0->NDTR = HALF;
    DMA1_Stream0->CR   = DMA_SxCR_DBM            /* two pointers, ping-pong */
                       | DMA_SxCR_MINC
                       | (1u << DMA_SxCR_MSIZE_Pos)   /* 16-bit memory */
                       | (1u << DMA_SxCR_PSIZE_Pos)   /* 16-bit peripheral */
                       | DMA_SxCR_TCIE
                       | DMA_SxCR_EN;
    /* The request number that connects ADC1 to this stream comes from the
       DMAMUX request table in RM0455. It is not the same number as RM0433. */
}

void DMA1_Stream0_IRQHandler(void)
{
    DMA1->LIFCR = DMA_LIFCR_CTCIF0;
    /* CT says which pointer the engine is filling NOW, so the other is ready */
    ready_ptr = (DMA1_Stream0->CR & DMA_SxCR_CT) ? buf_a : buf_b;
    ready_seq++;
    marker_toggle();
}
\end{ccode}
Two traps live in these twenty lines and both belong to this family. The first
is the cache: with the data cache enabled and the buffers in the main memory
region, the processor can read a stale line that the engine has already
overwritten. The fix is to invalidate the block before reading it, and the
complete treatment is chapter 19. The second is reach: the tightly coupled
memories are fast and the main transfer engines cannot see them at all, so a
buffer moved there for speed produces a transfer that appears to run and writes
nothing. That is item 16 on the confirm list and the most common fault in this
family.

\item \textbf{Run the witness.} On the Pi, scan one channel continuously and
write the samples with the metadata that makes them interpretable later. The
API names below follow the HAT library's published examples; check them against
the version you install rather than against this page.
\begin{pycode}
# witness/scan.py -- runs on the Raspberry Pi
import json, time
import numpy as np
from daqhats import mcc118, OptionFlags, hat_list, HatIDs

RATE = 100000.0          # requested; the library reports what it granted
SECONDS = 60
CHANNEL = 0

board = mcc118(hat_list(filter_by_id=HatIDs.MCC_118)[0].address)
actual = board.a_in_scan_actual_rate(1, RATE)
board.a_in_scan_start(1 << CHANNEL, int(actual * SECONDS), actual,
                      OptionFlags.CONTINUOUS)
t0 = time.time()
chunks = []
while sum(len(c) for c in chunks) < actual * SECONDS:
    r = board.a_in_scan_read_numpy(-1, 1.0)
    if r.hardware_overrun or r.buffer_overrun:
        raise SystemExit("witness overran; lower the rate or the duration")
    chunks.append(r.data)
board.a_in_scan_stop()
board.a_in_scan_cleanup()

v = np.concatenate(chunks)
np.save("samples.npy", v)
open("meta.json", "w").write(json.dumps(
    {"requested_rate": RATE, "actual_rate": actual,
     "started_unix": t0, "channel": CHANNEL, "samples": int(v.size)}))
\end{pycode}
Record the rate the library granted rather than the rate you asked for. A HAT
that quietly rounds the scan rate to something it can divide exactly turns a
clean measurement into a systematic error of the same size as the effect.

\item \textbf{Turn edges into a rate with an uncertainty.} The analysis is
short and every line of it is a decision worth defending in a README.
\begin{pycode}
# witness/rate.py
import json
import numpy as np

v = np.load("samples.npy")
meta = json.load(open("meta.json"))
fs = meta["actual_rate"]

# Threshold midway between the two settled levels, not at a fixed voltage:
# the marker's high level depends on the series resistor and the HAT's input.
lo, hi = np.percentile(v, 5), np.percentile(v, 95)
thr = 0.5 * (lo + hi)

above = v > thr
cross = np.flatnonzero(np.diff(above.astype(np.int8)) != 0)
# Linear interpolation between the two samples that straddle the threshold.
frac = (thr - v[cross]) / (v[cross + 1] - v[cross])
t_edge = (cross + frac) / fs

n = np.arange(t_edge.size)
slope, intercept = np.polyfit(n, t_edge, 1)
resid = t_edge - (slope * n + intercept)

period = slope                      # one edge per sampling interrupt
rate = 1.0 / period
# Standard error of the slope of a straight-line fit.
se = resid.std(ddof=2) / (n.std() * np.sqrt(n.size))
print(f"edges           {t_edge.size}")
print(f"sample period   {period*1e6:.4f} us")
print(f"sample rate     {rate:.4f} Hz  +/- {rate*se/period:.4f} Hz")
print(f"jitter, rms     {resid.std()*1e9:.1f} ns")
print(f"jitter, peak    {np.abs(resid).max()*1e9:.1f} ns")
\end{pycode}
Report all four numbers. The rate without the jitter hides a build that is on
average right and individually wrong, which is exactly what the tick build with
a long conversion wait looks like.

\item \textbf{Build the gated frequency counter on the board.} The witness is
on the other side of the room and needs a Pi. The board can measure a periodic
signal by itself, and from here on it does. One 32-bit timer counts edges
arriving on its external trigger input; a second timer, clocked from the system
clock, opens and closes the gate.
\begin{ccode}
/* instr/freqcount.c: counts edges on TIM5_ETR over a gate from TIM7. */
void freq_start(uint32_t gate_ms)
{
    TIM5->PSC  = 0u;
    TIM5->ARR  = 0xFFFFFFFFu;
    TIM5->SMCR = TIM_SMCR_ECE;        /* external clock mode 2, on ETR */
    TIM5->CR1  = TIM_CR1_CEN;

    TIM7->PSC = 28000u - 1u;          /* 280 MHz -> 10 kHz */
    TIM7->ARR = (gate_ms * 10u) - 1u;
    TIM7->DIER = TIM_DIER_UIE;
    TIM7->CR1  = TIM_CR1_CEN;
}

static volatile uint32_t last, delta;
void TIM7_IRQHandler(void)
{
    TIM7->SR = 0u;
    uint32_t now = TIM5->CNT;
    delta = now - last;               /* unsigned, so the wrap is free */
    last  = now;
}
\end{ccode}
The resolution of this instrument is one count in the gate. On a 1 kHz signal
with a one second gate that is one part in a thousand, which is far too coarse
to see the effect this chapter is after. The upgrade is the classic one and is
worth the extra twenty lines: stop counting edges in a fixed window and instead
capture the system-clock time of the Nth edge. With a 32-bit capture timer at
280 MHz, a thousand periods are timed to about 3.6 ns, which is under four parts
in a billion of the interval. Resolution then stops being the limit and the
oscillator becomes the limit, which is the honest place for it to be.

\item \textbf{Test the analysis before you trust it.} The rate script is the
instrument's most fragile part, so it is tested on synthetic input with a known
answer, on the host, in continuous integration. This is the same discipline
chapter 12 turns into a rig.
\begin{pycode}
# witness/test_rate.py
import numpy as np
from rate import edges_and_fit       # refactored out of the script

def test_known_rate():
    fs, f, secs = 100000.0, 500.0, 5.0
    t = np.arange(int(fs * secs)) / fs
    v = 1.65 + 1.5 * np.sign(np.sin(2 * np.pi * f * t))
    v += np.random.default_rng(0).normal(0, 0.01, v.size)
    period, jitter = edges_and_fit(v, fs)
    assert abs(1.0 / period - 2 * f) < 0.05     # one edge per half period
    assert jitter < 2e-6
\end{pycode}
\end{steps}

\subsection*{Build, flash and debug}

\diagram{c06_timing}{One millisecond as three of the four parties see it. The witness's sample instants are its own and are never aligned to the board's; the interpolated crossing between two of them is what makes a rate out of a square wave.}

Three builds, one application, three binaries. The CMake target names carry the
variant so a flashed board can be identified from its console banner.

\begin{shellcode}
cmake -B build -DCMAKE_TOOLCHAIN_FILE=cmake/arm-none-eabi.cmake
cmake --build build --target acq_systick acq_timer acq_dma_double -j
probe-rs run --chip STM32H7A3ZITx build/acq_dma_double.elf
\end{shellcode}

On the Pi, the witness runs and the analysis follows it, so a run produces a
number and not a folder of samples somebody will interpret later.

\begin{shellcode}
ssh pi@witness 'cd witness && python scan.py && python rate.py' | tee run.log
\end{shellcode}

\begin{note}{When the rate is right and the jitter is not}
The usual cause is that the marker is toggled after work rather than before it,
so the edge carries the duration of whatever the interrupt did first. Move the
toggle to the first statement and measure again. If the jitter survives that,
the next suspects in order are a higher-priority interrupt that preempts the
sampling one, a wait state on the memory the buffer lives in, and the converter
sampling time rather than the trigger. The gated counter cannot separate those;
the cycle counter placed around each candidate can.
\end{note}

If the board stops responding to the debugger after a build that enables the
transfer engine, the recovery is the same as chapter 1's: hold the reset while
the flashing tool connects, or use the probe's connect-under-reset option. No
option bytes are written in this chapter, so there is nothing here that can put
the board beyond that recovery.

\subsection*{Verification and acceptance criteria}

\begin{itemize}
\item The witness reports a sample rate whose stated uncertainty does not
contain the difference between the three builds, or it reports that it does and
the chapter says the three are indistinguishable at this rate. Either result is
publishable; a result with no uncertainty is not.
\item The three builds compile from one unmodified application source, proven by
a checksum of that file being identical in all three build logs.
\item The cycle counter's self-test runs on a board that has never had a
debugger attached in that power cycle, and the result is recorded in the README
either way. If the trace counter does not run without a debugger, the timer
backend is used and the chapter says so.
\item The overrun count is zero over a ten minute run for the double-buffered
build, and is greater than zero when the processing is deliberately slowed past
the block budget. A counter that never increments has not been tested.
\item The rate analysis passes its synthetic test on the host, in continuous
integration, with no board attached.
\item The marker toggle's cost in cycles is measured and written in the
repository, because seven later chapters inherit it.
\end{itemize}

\subsection*{Variants}

\begin{table}[H]\centering\small
\begin{tabular}{p{24mm}p{26mm}p{44mm}p{22mm}p{20mm}}
\toprule
Axis & Variant & What changes & Cost & Built in full in \\
\midrule
Time and safety & SysTick & The core's own exception at 1 kHz. No timer is spent, and the conversion wait sits inside the exception & Jitter tied to whatever else preempts it & Here \\
Time and safety & General timer & The timer's trigger output starts the conversion with no software in the path, which is the first build whose rate is a hardware property & One timer & Here \\
Execution model & DMA double buffered & Two memory pointers and a current-target bit, so the application always has a whole half to itself & Cache maintenance becomes mandatory & Here \\
Execution model & DMA circular with half transfer & The same ping-pong with one pointer and a half-transfer interrupt. Available on every transfer engine here & The idle half's address cannot be changed & Chapter 4 \\
Execution model & Polled loop or an interrupt with a flag & Read the converter in the main loop, or let the end-of-conversion interrupt raise a flag the loop acts on & No fixed rate at all, or one sample of latency & Chapter 3 \\
Execution model & RTOS task & A task blocks on a notification from the transfer interrupt & A kernel and its tick & Chapter 20 \\
Time and safety & Low-power timer & The same 1 kHz from a timer that survives the deeper sleep states & Fewer features & Chapter 7 \\
Peripheral & Hardware oversampling & The converter averages in hardware before the transfer, so the rate falls and the resolution rises with no processor involvement & Latency and a decision about the shift & Chapter 18 \\
Intelligence & Edge host & Stream every sample to the Pi and do the work there & A framed link and its bandwidth & Chapter 9 and 11 \\
\bottomrule
\end{tabular}
\caption{Variants for chapter 6. The three built here are the three whose difference the witness can actually resolve; the rest are referenced because their difference shows up at rates this chapter does not reach.}
\end{table}

The comparison across the top three rows is the part this book has not found
published. SOURCE.md records as a gap that nobody has measured several execution
models solving one problem on one board with a stated method. This chapter
measures three of them, states the method, and leaves the remaining three to
chapters 3, 4 and 20 using the same witness and the same analysis script, which
is what makes the numbers comparable across the volume.

\subsection*{Pitfalls}

\begin{itemize}
\item Copying a converter trigger source number or a transfer request number
from a project for the better-known member of this family. The tables differ.
The failure is a converter that never starts, which reads as a dead transfer.
\item Placing the sample buffer in tightly coupled memory because it is fast.
The main transfer engines cannot reach it. The transfer appears to run.
\item Leaving the data cache enabled and reading the buffer without
invalidating it. The processor returns bytes that were correct one block ago.
\item Believing the requested scan rate on the witness rather than the rate it
granted. A rounded scan rate is a systematic error, not noise.
\item Toggling the marker after the interrupt's work instead of before it, then
reporting the resulting spread as the board's jitter.
\item Comparing two clocks that share an oscillator. If the witness is ever
driven from the board, or the board from the witness, the measurement proves
nothing.
\item Reporting an average rate with no uncertainty and no jitter, or assuming
the trace cycle counter runs without a debugger. The second is item 25 on the
confirm list precisely because the answer decides which instrument this book
uses.
\end{itemize}

\subsection*{Best practices applied}

\begin{itemize}
\item Every measurement names its instrument. The rate and the jitter come from
the MCC 118; the cycle costs come from this chapter's own wrapper; the cells
that neither has filled say so.
\item The instrument is tested before it is trusted. The analysis script runs
against synthetic input with a known answer, and the cycle counter tests itself
at startup rather than assuming a debug block is enabled.
\item The three builds share one application source, so the comparison measures
the acquisition strategy and nothing else, and the uncertainty reported with
every number is derived from the fit rather than asserted.
\item The overrun path is exercised deliberately, because an error counter that
has never incremented is a decoration.
\item Facts that are not settled are written as confirm-list items with a
reference, not as assertions with a plausible number.
\end{itemize}

\subsection*{Stretch goals}

\begin{itemize}
\item Replace the gated counter with the reciprocal counter sketched in step 9
and show the resolution improving by three orders of magnitude while the stated
accuracy does not move, because the oscillator has become the limit.
\item Run the witness against the board's low-speed crystal instead of its
sampling interrupt, by routing a divided crystal output to the marker pin, and
compare the two oscillators the board carries. Chapter 7 needs that number.
\item Raise the rate to 10 kHz and then 50 kHz and find where each of the three
builds stops meeting its budget. The tick build will be the first to stop and
the witness can say by how much.
\item Publish the analysis script and its synthetic test as a small standalone
package, since no permissively licensed tool in the pool turns a sampled square
wave into a rate with an uncertainty.
\end{itemize}

\subsection*{Roadmap and next steps}

Chapter 7 takes the timer built here into the sleep states, where the
peripheral bus clock stops and the general timer stops with it, which is why the
low-power timer exists and why that chapter rather than this one builds it.

For going deeper, the community embedded roadmap is the map worth following,
and its central advice, that projects settle questions reading leaves open, is
this volume's thesis as well. The written series behind the five-step
progression cited above is the best companion to this chapter specifically, and
reading it against RM0455 rather than against the part it was written for is
itself a useful exercise. For the scheduling mathematics that decides when a
sampling deadline is actually met rather than usually met, the university
specialisation on real-time embedded systems listed in Appendix H is the only
free course in the pool that teaches it properly; nothing else here does. The
vendor's own free training modules cover this family, though there is no course
specific to this part.

\subsection*{Portfolio evidence}

\begin{itemize}
\item A repository with three build targets, one application source, and a
README that opens with the architecture figure and states the measured rate with
its uncertainty in the first paragraph.
\item The run log from the witness, showing the granted scan rate, the edge
count, the rate, the standard error and both jitter figures.
\item A short written method, dated and written before the first run, saying
what would count as a negative result.
\item The cycle counter wrapper and the frequency counter as a small library
with their self-tests, reused unchanged by later chapters, which is better
evidence than a larger piece of code used once.
\end{itemize}

\subsection*{Sources}

Normative references:
\begin{itemize}
\item Reference manual RM0455, for the timer inventory, the trigger output
selection, the converter's external trigger table and the transfer request
table. Not RM0433.
\item The STM32H7A3xI datasheet, for the alternate-function table that decides
which pins carry the external trigger input and the capture inputs, and for the
converter's channel mapping.
\item The board user manual for MB1363, for which Zio pins remain free and
carry a timer channel.
\item The Cortex-M7 technical reference manual and the architecture reference
manual, for the trace block, the cycle counter and the conditions under which it
counts.
\item The MCC 118 user guide, for the single-channel and aggregate scan rates,
the input range and whether an external trigger exists.
\end{itemize}

Reusable implementations:
\begin{itemize}
\item Ian Ross's five-step timer and converter progression on a Cortex-M7
Nucleo, MIT, with the written series that explains each step.\newline
\url{https://github.com/ian-ross/stm32-timer-adc-dma}\newline
\url{https://www.skybluetrades.net/blog/2020/11/2020-11-26-stm32-timer-adc-dma-3/}
\item Majerle's transfer-driven serial examples, MIT, for the half and full
transfer pattern reused here.\newline
\url{https://github.com/MaJerle/stm32-usart-uart-dma-rx-tx}
\item GorgonMeducer's performance counter, Apache-2.0, which chooses the system
tick over the trace counter for exactly the reason step 2 tests.\newline
\url{https://github.com/GorgonMeducer/perf_counter}
\item The flashing and debugging tool whose target list includes this die.\newline
\url{https://github.com/probe-rs/probe-rs}
\end{itemize}
